\clearpage
\section{Eksperimen}

Bagian ini berisi hasil eksekusi program terhadap seluruh kasus uji pada Bab 5 spesifikasi. Berkas masukannya tersimpan di folder \texttt{test/} dengan pola nama \texttt{spl\_*.txt}, \texttt{det\_*.txt}, \texttt{inverse\_*.txt}, \texttt{interpolasi\_*.txt}, dan \texttt{regresi\_*.txt}. Angka pada bagian \textit{Output} adalah keluaran pustaka kelompok yang sudah dibulatkan ke 3 angka di belakang koma, sama seperti yang ditampilkan program. Sebagai pembanding, kami menghitung solusi eksak dengan aritmetika pecahan (\textit{rational}) di SymPy.

\subsection{Determinan Matriks}

Setiap kasus diuji dengan metode reduksi baris (OBE) dan ekspansi kofaktor.

\subsubsection{Kasus Uji 1}
\begin{kasus}
\ubag{Input}
\[
\begin{bmatrix}
0 & 2 & -1 & 3 \\
4 & 1 & 2 & -2 \\
-2 & 3 & 0 & 1 \\
1 & -1 & 5 & 2
\end{bmatrix}
\]
\ubag{Output}
\begin{center}
OBE: $\det(A) = -259$ \qquad Kofaktor: $\det(A) = -259$
\end{center}
\begin{center}
    \begin{minipage}{0.48\textwidth}
        \centering
        \tangkapan{det_1_obe}
    \end{minipage}\hfill
    \begin{minipage}{0.48\textwidth}
        \centering
        \tangkapan{det_1_kof}
    \end{minipage}
\end{center}
\end{kasus}

\subsubsection{Kasus Uji 2}
\begin{kasus}
\ubag{Input}
\[
\begin{bmatrix}
2 & -1 & 3 & 0 & 4 \\
1 & 2 & -2 & 5 & -1 \\
3 & 1 & 1 & 2 & 0 \\
4 & -1 & 3 & 2 & 0 \\
5 & 0 & 4 & 2 & 4
\end{bmatrix}
\]
\ubag{Output}
\begin{center}
OBE: $\det(A) = 0$ \qquad Kofaktor: $\det(A) = 0$
\end{center}
\begin{center}
    \begin{minipage}{0.48\textwidth}
        \centering
        \tangkapan{det_2_obe}
    \end{minipage}\hfill
    \begin{minipage}{0.48\textwidth}
        \centering
        \tangkapan{det_2_kof}
    \end{minipage}
\end{center}
\end{kasus}
\subsubsection{Kasus Uji 3}

\begin{kasus}
\ubag{Input}
\[
\begin{bmatrix}
1/2 & 1 & 0 & 1/2 \\
1/2 & 5/4 & 1/4 & 1/2 \\
-1/10 & -2/5 & 1/5 & -3/10 \\
3/5 & 5/4 & 9/20 & 3/5
\end{bmatrix}
\]
Program tidak menerima pecahan, sehingga berkas masukan
(\texttt{det\_5\_1\_3.txt}) ditulis dalam desimal ($0.5$, $1.25$, $0.25$,
$-0.1$, $-0.4$, $0.2$, $-0.3$, $0.6$, $0.45$).
\ubag{Output}
\begin{center}
OBE: $\det(A) = 0.01$ \qquad Kofaktor: $\det(A) = 0.01$ \qquad
Eksak: $1/100 = 0.01$
\end{center}
\begin{center}
    \begin{minipage}{0.48\textwidth}
        \centering
        \tangkapan{det_3_obe}
    \end{minipage}\hfill
    \begin{minipage}{0.48\textwidth}
        \centering
        \tangkapan{det_3_kof}
    \end{minipage}
\end{center}
\end{kasus}

\textbf{Analisis.} Kasus 1, 2, dan 3 semuanya menghasilkan nilai yang sama persis untuk kedua metode. Berdasarkan pengujian matriks 5x5, kedua metode sama-sama menghasilkan determinan 0 lalu algoritma Kofaktor bekerja  dengan menjadikan Baris 1 sebagai acuan untuk meng-skip minor bernilai nol sementara OBE mendeteksi elemen pivot nol di langkah ke-14 yang otomatis membuat perkalian diagonal utamanya mutlak bernilai 0.
 Keakuratan serupa juga terbukti pada matriks desimal pecahan 4x4 yang menghasilkan determinan 0.01. Kofaktor  mengoptimasi komputasi dengan memotong perhitungan pada elemen nol di Baris 1 sedangkan OBE berhasil mereduksi matriks menjadi segitiga atas dengan penyesuaian tanda kali -1 akibat operasi tukar baris sehingga hasil kali seluruh elemen diagonal utamanya secara konsisten berujung pada nilai 0.01 yang identik.

\subsection{Invers Matriks}

Setiap kasus diuji dengan metode augmentasi dan metode adjoin.

\subsubsection{Kasus Uji 1}
\begin{kasus}
\ubag{Input}
\[
\begin{bmatrix}
0 & 1 & 2 & -1 \\
1 & 0 & 1 & 2 \\
2 & 1 & 0 & 1 \\
-1 & 2 & 1 & 0
\end{bmatrix}
\]
\ubag{Output}
Augmentasi dan adjoin menghasilkan matriks yang sama:
\[
A^{-1} = \begin{bmatrix}
0.2 & -0.1 & 0.4 & -0.3 \\
-0.1 & -0.2 & 0.3 & 0.4 \\
0.4 & 0.3 & -0.2 & -0.1 \\
-0.3 & 0.4 & -0.1 & 0.2
\end{bmatrix}
\]
\textbf{Augmentasi}
\tangkapan{hasil_inverse_1_aug_1}
\tangkapan{hasil_inverse_1_aug_2}
\tangkapan{hasil_inverse_1_aug_3}
\tangkapan{hasil_inverse_1_aug_4}

\textbf{Kofaktor}
\tangkapan{hasil_inverse_1_kof_1}
\tangkapan{hasil_inverse_1_kof_2}
\tangkapan{hasil_inverse_1_kof_3}
\end{kasus}

\subsubsection{Kasus Uji 2}
\begin{kasus}
\ubag{Input}
\[
\begin{bmatrix}
1 & 2 & 0 & -1 & 3 \\
2 & -1 & 4 & 0 & 1 \\
0 & 3 & 1 & 2 & -2 \\
4 & 0 & 5 & 1 & 2 \\
3 & 1 & 4 & -1 & 4
\end{bmatrix}
\]
\ubag{Output}
Kedua metode: \texttt{matriks tidak memiliki balikan (Matriks singular, determinan = 0)}.

\textbf{Augmentasi}
\tangkapan{inverse_2_aug_1}
\tangkapan{inverse_2_aug_2}
\tangkapan{inverse_2_aug_3}
\tangkapan{inverse_2_aug_4}
\tangkapan{inverse_2_aug_5}
\tangkapan{inverse_2_aug_6}

\textbf{Kofaktor}
\tangkapan{inverse_2_kof_1}
\end{kasus}

\textbf{Analisis.} Pada kasus 1, kedua metode menghasilkan invers matriks yang sama. Hasil tersebut juga sesuai dengan nilai yang ditulis dalam bentuk nilai sebenarnya (tanpa pembulatan desimal), yaitu matriks dengan elemen-elemen yang merupakan kelipatan $1/10$. Selain itu, karena matriks $A$ simetris, maka matriks balikannya pun juga simetris.

Sementara, di kasus 2, matriks tidak memiliki invers matriks karena baris kelima merupakan hasil penjumlahan baris pertama dan baris kedua. Karena itu, baris-baris matriks tidak bebas secara linier dan $\det(A) = 0$. Pada metode augmentasi, kondisi ini baru terlihat ketika tidak ditemukan pivot yang dapat digunakan pada salah satu kolom. Sementara itu, metode adjoin terlebih dahulu memeriksa determinan. Karena determinannya nol, jadi proses mencari balikan dihentikan.

\subsection{Sistem Persamaan Linier $Ax = b$}

Setiap kasus dijalankan dengan keempat metode. Metode matriks balikan dan kaidah Cramer hanya bisa dipakai jika $A$ persegi dan $\det(A) \neq 0$ yang dimana jika syarat itu tidak terpenuhi, program menampilkan pesan alasannya.

\subsubsection{Kasus Uji 1}
\begin{kasus}
\ubag{Input}
\[
A = \begin{bmatrix}
0 & 2 & -1 & 1 \\
1 & -1 & 2 & 3 \\
2 & 1 & 1 & -1 \\
-1 & 3 & 0 & 2
\end{bmatrix}, \quad
b = \begin{bmatrix} -5 \\ 15 \\ 1 \\ -3 \end{bmatrix}
\]
\ubag{Output}
\begin{center}
Gauss, Gauss-Jordan, Balikan, Cramer: $x_1 = 1$, $x_2 = -2$, $x_3 = 3$, $x_4 = 2$
\end{center}
\textbf{Eliminasi Gauss}
\tangkapan{spl_1_gauss_1}
\tangkapan{spl_1_gauss_2}
\tangkapan{spl_1_gauss_3}

\textbf{Eliminasi Gauss-Jordan}
\tangkapan{spl_1_gj_1}
\tangkapan{spl_1_gj_2}
\tangkapan{spl_1_gj_3}

\textbf{Metode Matriks Balikan}
\tangkapan{spl_1_balikan_1}
\tangkapan{spl_1_balikan_2}
\tangkapan{spl_1_balikan_3}

\textbf{Kaidah Cramer}
\tangkapan{spl_1_cramer_1}
\end{kasus}

\subsubsection{Kasus Uji 2}
\begin{kasus}
\ubag{Input}
\[
A = \begin{bmatrix}
1 & 2 & -1 & 0 & 3 \\
2 & 4 & 1 & -1 & 5 \\
-1 & -2 & 2 & 3 & -1
\end{bmatrix}, \quad
b = \begin{bmatrix} 2 \\ 10 \\ -1 \end{bmatrix}
\]
\ubag{Output}
Gauss dan Gauss-Jordan (tak hingga solusi):
\begin{align*}
x_1 &= 3.9 - 2t_1 - 2.9t_2, & x_2 &= t_1, & x_3 &= 1.9 + 0.1t_2, \\
x_4 &= -0.3 - 0.7t_2, & x_5 &= t_2 .
\end{align*}
Balikan dan Cramer: \texttt{tidak bisa dipakai karena jumlah persamaan tidak sama dengan jumlah variabel}.
\textbf{Eliminasi Gauss}
\tangkapan{spl_2_gauss_1}
\tangkapan{spl_2_gauss_2}

\textbf{Eliminasi Gauss-Jordan}
\tangkapan{spl_2_gj_1}
\tangkapan{spl_2_gj_2}

\textbf{Metode Matriks Balikan}
\tangkapan{spl_2_balikan_1}

\textbf{Kaidah Cramer}
\tangkapan{spl_2_cramer_1}
\end{kasus}

\subsubsection{Kasus Uji 3}
\begin{kasus}
\ubag{Input}
\[
A = \begin{bmatrix}
1 & -2 & 3 & 1 \\
2 & 1 & -1 & 4 \\
0 & 3 & 2 & -2 \\
3 & -1 & 2 & 5
\end{bmatrix}, \quad
b = \begin{bmatrix} 4 \\ -1 \\ 3 \\ 4 \end{bmatrix}
\]
\ubag{Output}
Gauss dan Gauss-Jordan: \texttt{solusi tidak ada}.\\
Balikan dan Cramer: \texttt{tidak bisa dipakai karena det(A) = 0}.
\textbf{Eliminasi Gauss}
\tangkapan{spl_3_gauss_1}
\tangkapan{spl_3_gauss_2}

\textbf{Eliminasi Gauss-Jordan}
\tangkapan{spl_3_gj_1}
\tangkapan{spl_3_gj_2}
\tangkapan{spl_3_gj_3}

\textbf{Metode Matriks Balikan}
\tangkapan{spl_3_balikan_1}

\textbf{Kaidah Cramer}
\tangkapan{spl_3_cramer_1}
\end{kasus}

\subsubsection{Kasus Uji 4: Matriks \textit{Discrete Laplacian}}
\begin{kasus}
\ubag{Input}
$A_n$ bernilai 2 pada diagonal, $-1$ tepat di atas dan di bawah diagonal, dan 0 di tempat lain dan $b$ berisi angka 1. Diuji untuk $n = 6$ dan $n = 10$.
\ubag{Output}
\begin{center}
\begin{tabular}{lcccccccccc}
\toprule
$n$ & $x_1$ & $x_2$ & $x_3$ & $x_4$ & $x_5$ & $x_6$ & $x_7$ & $x_8$ & $x_9$ & $x_{10}$ \\
\midrule
6  & 3 & 5 & 6 & 6 & 5 & 3 & & & & \\
10 & 5 & 9 & 12 & 14 & 15 & 15 & 14 & 12 & 9 & 5 \\
\bottomrule
\end{tabular}
\end{center}
Keempat metode memberi hasil yang sama setelah pembulatan.
\textbf{\boldmath$n = 6$}

\textbf{Eliminasi Gauss}
\tangkapan{spl_4_n6_gauss_1}
\tangkapan{spl_4_n6_gauss_2}
\tangkapan{spl_4_n6_gauss_3}

\textbf{Eliminasi Gauss-Jordan}
\tangkapan{spl_4_n6_gj_1}
\tangkapan{spl_4_n6_gj_2}
\tangkapan{spl_4_n6_gj_3}
\tangkapan{spl_4_n6_gj_4}
\tangkapan{spl_4_n6_gj_5}

\textbf{Metode Matriks Balikan}
\tangkapan{spl_4_n6_balikan_1}
\tangkapan{spl_4_n6_balikan_2}
\tangkapan{spl_4_n6_balikan_3}
\tangkapan{spl_4_n6_balikan_4}
\tangkapan{spl_4_n6_balikan_5}
\tangkapan{spl_4_n6_balikan_6}

\textbf{Kaidah Cramer}
\tangkapan{spl_4_n6_cramer_1}
\tangkapan{spl_4_n6_cramer_2}

\textbf{\boldmath$n = 10$}

\textbf{Eliminasi Gauss}
\tangkapan{spl_4_n10_gauss_1}
\tangkapan{spl_4_n10_gauss_2}
\tangkapan{spl_4_n10_gauss_3}
\tangkapan{spl_4_n10_gauss_4}
\tangkapan{spl_4_n10_gauss_5}
\tangkapan{spl_4_n10_gauss_6}

\textbf{Eliminasi Gauss-Jordan}
\tangkapan{spl_4_n10_gj_1}
\tangkapan{spl_4_n10_gj_2}
\tangkapan{spl_4_n10_gj_3}
\tangkapan{spl_4_n10_gj_4}
\tangkapan{spl_4_n10_gj_5}
\tangkapan{spl_4_n10_gj_6}
\tangkapan{spl_4_n10_gj_7}
\tangkapan{spl_4_n10_gj_8}
\tangkapan{spl_4_n10_gj_9}
\tangkapan{spl_4_n10_gj_10}
\tangkapan{spl_4_n10_gj_11}
\tangkapan{spl_4_n10_gj_12}
\tangkapan{spl_4_n10_gj_13}
\tangkapan{spl_4_n10_gj_14}
\tangkapan{spl_4_n10_gj_15}
\tangkapan{spl_4_n10_gj_16}
\tangkapan{spl_4_n10_gj_17}

\textbf{Metode Matriks Balikan}
\tangkapan{spl_4_n10_balikan_1}
\tangkapan{spl_4_n10_balikan_2}
\tangkapan{spl_4_n10_balikan_3}
\tangkapan{spl_4_n10_balikan_4}
\tangkapan{spl_4_n10_balikan_5}
\tangkapan{spl_4_n10_balikan_6}
\tangkapan{spl_4_n10_balikan_7}
\tangkapan{spl_4_n10_balikan_8}
\tangkapan{spl_4_n10_balikan_9}
\tangkapan{spl_4_n10_balikan_10}
\tangkapan{spl_4_n10_balikan_11}
\tangkapan{spl_4_n10_balikan_12}
\tangkapan{spl_4_n10_balikan_13}
\tangkapan{spl_4_n10_balikan_14}
\tangkapan{spl_4_n10_balikan_15}
\tangkapan{spl_4_n10_balikan_16}
\tangkapan{spl_4_n10_balikan_17}
\tangkapan{spl_4_n10_balikan_18}

\textbf{Kaidah Cramer}
\tangkapan{spl_4_n10_cramer_1}
\tangkapan{spl_4_n10_cramer_2}
\tangkapan{spl_4_n10_cramer_3}
\tangkapan{spl_4_n10_cramer_4}
\end{kasus}

\textbf{Analisis.} Pada kasus 1, keempat metode memberikan jawaban yang sama dengan solusi eksak. Jika dilihat sebelum dibulatkan, nilainya sedikit berbeda antarmetode. Contohnya, $x_1$ bernilai $1.0000000000000002$ pada eliminasi Gauss dan $1.0000000000000004$ pada Gauss-Jordan. Perbedaan ini baru muncul di digit ke-16, sesuai dengan batas ketelitian \texttt{double}, dan hilang setelah dibulatkan ke 3 angka di belakang koma.

Kasus 2 hanya memiliki 3 persamaan untuk 5 variabel, jadi paling sedikit ada 2 variabel bebas. Program memilih $x_2$ dan $x_5$ sebagai parameter karena kolom keduanya tidak mendapat pivot. Kami membandingkannya dengan hasil SymPy dan bentuknya sama: $x_1 = 39/10 - 2x_2 - 29x_5/10$, $x_3 = 19/10 + x_5/10$, dan $x_4 = -3/10 - 7x_5/10$. Metode balikan dan Cramer langsung ditolak karena matriks $A$ tidak persegi.

Pada kasus 3, $\operatorname{rank}(A) = 3$, sedangkan $\operatorname{rank}([A\,|\,b]) = 4$. Akibatnya, setelah eliminasi muncul baris $0 = c$ dengan $c \neq 0$ sehingga SPL tidak memiliki solusi. Metode balikan dan Cramer sudah berhenti lebih awal karena $\det(A) = 0$, tetapi pesan itu belum menjelaskan apakah SPL-nya tidak punya solusi atau punya banyak solusi. Untuk kasus seperti ini, hanya eliminasi Gauss dan Gauss-Jordan yang memberikan jawaban lengkap.

Kasus 4 kami pakai untuk melihat kestabilan numerik keempat metode. Solusi eksaknya mengikuti rumus $x_i = i(n+1-i)/2$, dan hasil program untuk $n = 6$ maupun $n = 10$ cocok dengan rumus tersebut. Matriks Laplacian bersifat dominan diagonal sehingga \textit{partial pivoting} tidak perlu menukar baris serta pada tangkapan layar, jumlah tukar baris tercatat 0. Galat terbesar sebelum pembulatan muncul pada kaidah Cramer untuk $n = 10$, misalnya $x_2 = 9.000000000000004$. Hal ini masuk akal karena Cramer menghitung 11 determinan secara terpisah dan setiap determinan membawa galat pembulatannya sendiri.

\subsection{SPL Berbentuk Matriks \textit{Augmented}}

\subsubsection{Kasus Uji 1}
\begin{kasus}
\ubag{Input}
\[
\left[\begin{array}{cccc|c}
0 & 1 & -1 & 2 & -1 \\
2 & 0 & 1 & -1 & 3 \\
-1 & 2 & 0 & 1 & -3 \\
1 & -1 & 2 & 0 & 9
\end{array}\right]
\]
\ubag{Output}
\begin{center}
Keempat metode: $x_1 = 1$, $x_2 = -2$, $x_3 = 3$, $x_4 = 2$
\end{center}
\textbf{Eliminasi Gauss}
\tangkapan{spl_aug_1_gauss_1}
\tangkapan{spl_aug_1_gauss_2}
\tangkapan{spl_aug_1_gauss_3}

\textbf{Eliminasi Gauss-Jordan}
\tangkapan{spl_aug_1_gj_1}
\tangkapan{spl_aug_1_gj_2}
\tangkapan{spl_aug_1_gj_3}

\textbf{Metode Matriks Balikan}
\tangkapan{spl_aug_1_balikan_1}
\tangkapan{spl_aug_1_balikan_2}
\tangkapan{spl_aug_1_balikan_3}

\textbf{Kaidah Cramer}
\tangkapan{spl_aug_1_cramer_1}
\end{kasus}

\subsubsection{Kasus Uji 2}
\begin{kasus}
\ubag{Input}
\[
\left[\begin{array}{cccccc|c}
1 & 0 & 2 & -1 & 0 & 3 & -7 \\
0 & 1 & -1 & 2 & 1 & 0 & 6 \\
2 & -1 & 0 & 1 & 3 & -2 & 13 \\
1 & 1 & 1 & 1 & 1 & 1 & 3
\end{array}\right]
\]
\ubag{Output}
Gauss dan Gauss-Jordan (tak hingga solusi):
\begin{align*}
x_1 &= 5.8 - t_1 - 1.6t_2, & x_2 &= 2.6 - t_1 - 0.2t_2, & x_3 &= -3.4 + t_1 + 0.8t_2, \\
x_4 &= t_1, & x_5 &= t_2, & x_6 &= -2 .
\end{align*}
Balikan dan Cramer: tidak dapat dipakai karena matriks koefisien berukuran $4 \times 6$.
\textbf{Eliminasi Gauss}
\tangkapan{spl_aug_2_gauss_1}
\tangkapan{spl_aug_2_gauss_2}

\textbf{Eliminasi Gauss-Jordan}
\tangkapan{spl_aug_2_gj_1}
\tangkapan{spl_aug_2_gj_2}
\tangkapan{spl_aug_2_gj_3}

\textbf{Metode Matriks Balikan}
\tangkapan{spl_aug_2_balikan_1}

\textbf{Kaidah Cramer}
\tangkapan{spl_aug_2_cramer_1}
\end{kasus}

\textbf{Analisis.} Elemen $a_{11}$ pada kasus 1 bernilai 0. Tanpa pertukaran baris, langkah pertama eliminasi akan membagi dengan nol. Dengan \textit{partial pivoting}, program menukar baris 1 dan baris 2 karena nilai mutlak terbesar pada kolom pertama adalah 2, lalu eliminasi berjalan seperti biasa. Solusinya, $(1, -2, 3, 2)$, kebetulan sama dengan solusi SPL kasus 1 pada subbab sebelumnya.

Matriks koefisien pada kasus 2 berukuran $4 \times 6$ sehingga balikan dan Cramer tidak bisa dipakai. Hasil eliminasi memberikan dua variabel bebas, yaitu $x_4$ dan $x_5$. Yang perlu diperhatikan, $x_6$ tetap bernilai $-2$ berapa pun nilai parameternya, artinya setiap solusi SPL ini memenuhi $x_6 = -2$. Hasil ini cocok dengan perhitungan eksak.

\subsection{SPL Berbentuk Umum}

\subsubsection{Kasus Uji 1}
\begin{kasus}
\ubag{Input}
\[
\begin{cases}
3x_1 + 2x_2 - x_3 + 4x_4 = 5 \\
x_1 - 3x_2 + 2x_3 + x_4 = -2 \\
4x_1 + x_2 + 5x_3 - 2x_4 = 7 \\
2x_1 - x_2 + 3x_3 + x_4 = 3
\end{cases}
\]
\ubag{Output}
\begin{center}
\begin{tabular}{lcccc}
\toprule
 & $x_1$ & $x_2$ & $x_3$ & $x_4$ \\
\midrule
Program (4 metode) & $-0.114$ & $1.841$ & $1.432$ & $0.773$ \\
Nilai eksak & $-5/44$ & $81/44$ & $63/44$ & $17/22$ \\
\bottomrule
\end{tabular}
\end{center}
\textbf{Eliminasi Gauss}
\tangkapan{spl_umum_1_gauss_1}
\tangkapan{spl_umum_1_gauss_2}

\textbf{Eliminasi Gauss-Jordan}
\tangkapan{spl_umum_1_gj_1}
\tangkapan{spl_umum_1_gj_2}
\tangkapan{spl_umum_1_gj_3}

\textbf{Metode Matriks Balikan}
\tangkapan{spl_umum_1_balikan_1}
\tangkapan{spl_umum_1_balikan_2}
\tangkapan{spl_umum_1_balikan_3}

\textbf{Kaidah Cramer}
\tangkapan{spl_umum_1_cramer_1}
\end{kasus}

\subsubsection{Kasus Uji 2}
\begin{kasus}
\ubag{Input}
Sistem 9 persamaan dan 9 variabel pada spesifikasi 5.5.2, ditulis sebagai matriks \textit{augmented} $9 \times 10$ dengan mengelompokkan koefisien setiap variabel. Contohnya, persamaan keempat menjadi $0.03521x_1 + 0.68x_2 + 0.03521x_4 + 0.68x_6 + 0.03521x_7 + 0.55432x_9 = 8.47$.
\ubag{Output}
Gauss dan Gauss-Jordan (tak hingga solusi, satu parameter):
\begin{align*}
x_1 &= 5.776 - 0.462t, & x_2 &= 4.464 - 0.387t, & x_3 &= -0.24 + 0.849t, \\
x_4 &= -0.736 + 0.971t, & x_5 &= 5.224 - 0.538t, & x_6 &= 7.512 - 0.433t, \\
x_7 &= 4.227 - 0.416t, & x_8 &= 4.773 - 0.584t, & x_9 &= t .
\end{align*}
Balikan dan Cramer: \texttt{tidak bisa dipakai karena det(A) = 0}.
\textbf{Eliminasi Gauss}
\tangkapan{spl_umum_2_gauss_1}
\tangkapan{spl_umum_2_gauss_2}
\tangkapan{spl_umum_2_gauss_3}
\tangkapan{spl_umum_2_gauss_4}
\tangkapan{spl_umum_2_gauss_5}
\tangkapan{spl_umum_2_gauss_6}
\tangkapan{spl_umum_2_gauss_7}
\tangkapan{spl_umum_2_gauss_8}
\tangkapan{spl_umum_2_gauss_9}
\tangkapan{spl_umum_2_gauss_10}

\textbf{Eliminasi Gauss-Jordan}
\tangkapan{spl_umum_2_gj_1}
\tangkapan{spl_umum_2_gj_2}
\tangkapan{spl_umum_2_gj_3}
\tangkapan{spl_umum_2_gj_4}
\tangkapan{spl_umum_2_gj_5}
\tangkapan{spl_umum_2_gj_6}
\tangkapan{spl_umum_2_gj_7}
\tangkapan{spl_umum_2_gj_8}
\tangkapan{spl_umum_2_gj_9}
\tangkapan{spl_umum_2_gj_10}
\tangkapan{spl_umum_2_gj_11}
\tangkapan{spl_umum_2_gj_12}
\tangkapan{spl_umum_2_gj_13}
\tangkapan{spl_umum_2_gj_14}
\tangkapan{spl_umum_2_gj_15}

\textbf{Metode Matriks Balikan}
\tangkapan{spl_umum_2_balikan_1}

\textbf{Kaidah Cramer}
\tangkapan{spl_umum_2_cramer_1}
\end{kasus}

\textbf{Analisis.} Kasus 1 memiliki solusi tunggal dan keempat metode memberikan hasil yang sama. Nilai eksaknya berupa pecahan berpenyebut 44 (misalnya $x_1 = -5/44$) sehingga program hanya bisa menampilkannya sebagai desimal yang dibulatkan.

Kasus 2 awalnya terlihat seperti SPL persegi biasa dengan 9 persamaan dan 9 variabel. Namun, perhitungan eksak dengan pecahan menunjukkan $\operatorname{rank}(A) = \operatorname{rank}([A\,|\,b]) = 8$ dan $\det(A) = 0$, jadi jawaban solusi tak terhingga memang benar. Penyebabnya cukup mudah dilihat. Jika persamaan 1, 2, dan 3 dijumlahkan, ruas kirinya menjadi $x_1 + x_2 + \dots + x_9$, dan hal yang sama terjadi jika persamaan 7, 8, dan 9 dijumlahkan. Ruas kanannya pun sama, yaitu $10 + 12 + 9 = 31$ dan $11 + 8.5 + 11.5 = 31$. Dengan kata lain, salah satu persamaan merupakan kombinasi linier dari persamaan lain, tetapi SPL-nya tetap konsisten.

Kasus ini juga menguji toleransi nol pada program. Dalam \textit{floating point}, eliminasi tidak selalu menghasilkan nol yang tepat tetapi yang tersisa bisa berupa galat pembulatan yang sangat kecil. Jika batas nol terlalu ketat, sisa tersebut dapat dianggap sebagai pivot yang sah dan program akan melaporkan solusi tunggal yang keliru. Dengan toleransi relatif dari kelas \texttt{Toleransi}, kolom $x_9$ dikenali tidak memiliki pivot sehingga $x_9$ menjadi parameter $t$.

\subsection{Aplikasi Sistem Persamaan Linier}

\begin{kasus}
\ubag{Input}
Data luas rumah dan harga: $(1, 3)$, $(2, 5)$, $(3, 6)$. Dengan $\sum x_i = 6$, $\sum x_i^2 = 14$, $\sum y_i = 14$, dan $\sum x_i y_i = 31$, persamaan normalnya
\[
\left[\begin{array}{cc|c}
3 & 6 & 14 \\
6 & 14 & 31
\end{array}\right].
\]
\ubag{Output}
\begin{center}
Keempat metode: $\theta_0 = 1.667$, $\theta_1 = 1.5$, sehingga $y = 1.667 + 1.5x$.
\end{center}
\textbf{Eliminasi Gauss}
\tangkapan{spl_aplikasi_gauss_1}

\textbf{Eliminasi Gauss-Jordan}
\tangkapan{spl_aplikasi_gj_1}

\textbf{Metode Matriks Balikan}
\tangkapan{spl_aplikasi_balikan_1}

\textbf{Kaidah Cramer}
\tangkapan{spl_aplikasi_cramer_1}
\end{kasus}

\textbf{Analisis.} Nilai eksaknya adalah $\theta_0 = 5/3$ dan $\theta_1 = 3/2$. Garis $y = 1.667 + 1.5x$ tidak melewati ketiga titik secara tepat dan pada $x = 2$, model memprediksi $4.667$, sedangkan datanya bernilai 5. Hal ini wajar karena ketiga titik memang tidak segaris. Persamaan normal tidak mencari garis yang melewati semua titik, tetapi garis dengan jumlah kuadrat galat paling kecil. Ide yang sama kami pakai lagi pada regresi \textit{spline}, hanya saja jumlah basisnya lebih banyak.

\subsection{Studi Kasus Interpolasi Polinomial}

\subsubsection{Studi Kasus 1}
\begin{kasus}
\ubag{Input}
\begin{center}
\begin{tabular}{c|ccccccc}
$x$ & 0.2 & 0.4 & 0.6 & 0.8 & 1.0 & 1.2 & 1.4 \\ \hline
$f(x)$ & 0.182 & 0.336 & 0.470 & 0.588 & 0.693 & 0.788 & 0.875
\end{tabular}
\end{center}
Evaluasi di $x = 0.3$, $0.55$, $0.95$, dan $1.38$.
\ubag{Output}
\[
P(x) = 1.003x - 0.546x^2 + 0.479x^3 - 0.408x^4 + 0.208x^5 - 0.043x^6
\]
\begin{center}
\begin{tabular}{cccc}
\toprule
$x$ & $P(x)$ program & $\ln(1+x)$ & Selisih \\
\midrule
0.3  & 0.262 & 0.262 & $5.1 \times 10^{-4}$ \\
0.55 & 0.438 & 0.438 & $1.3 \times 10^{-4}$ \\
0.95 & 0.668 & 0.668 & $2.7 \times 10^{-5}$ \\
1.38 & 0.867 & 0.867 & $4.3 \times 10^{-4}$ \\
\bottomrule
\end{tabular}
\end{center}
\begin{center}
    \begin{minipage}{0.48\textwidth}
        \centering
        \tangkapan{interpol1}
    \end{minipage}\hfill
    \begin{minipage}{0.48\textwidth}
        \centering
        \tangkapan{interpol1_eval}
    \end{minipage}
\end{center}\end{kasus}

\textbf{Analisis.} Koefisien konstanta yang dihitung program bernilai $6.8 \times 10^{-13}$, sehingga tampil sebagai 0. Nilai pada tabel data cocok dengan $\ln(1+x)$ yang dibulatkan ke 3 desimal (misalnya $\ln 1.2 = 0.18232$), sehingga kami memakai $\ln(1+x)$ sebagai pembanding meskipun spesifikasi tidak menyebutkan fungsi asalnya. Selisihnya selalu di bawah $10^{-3}$ dan sebagian besar berasal dari pembulatan data masukan. Selisih terbesar muncul di $x = 0.3$ dan $x = 1.38$, di dekat ujung domain, sesuai dengan teori bahwa polinomial derajat tinggi paling tidak akurat di dekat ujung interval.

\subsubsection{Studi Kasus 2: Suhu Komponen}
\begin{kasus}
\ubag{Input}
\begin{center}
\begin{tabular}{c|cccccc}
$t$ (menit) & 0 & 2 & 4 & 6 & 8 & 10 \\ \hline
$T$ ($^\circ$C) & 25.0 & 31.2 & 38.6 & 47.5 & 58.1 & 70.8
\end{tabular}
\end{center}
Evaluasi di $t = 1$, $3.5$, $7$, dan $9.2$, serta cari waktu saat $T = 65^\circ$C.
\ubag{Output}
\[
T(t) = 25 + 2.8925t + 0.069792t^2 + 0.020313t^3 - 0.0018229t^4 + 0.000078125t^5
\]
\begin{center}
\begin{tabular}{cccc}
\toprule
$t = 1$ & $t = 3.5$ & $t = 7$ & $t = 9.2$ \\
\midrule
$27.981$ & $36.617$ & $52.571$ & $65.425$ \\
\bottomrule
\end{tabular}
\end{center}
Komponen mencapai $65^\circ$C pada $t \approx 9.14$ menit.
\begin{center}
    \begin{minipage}{0.48\textwidth}
        \centering
        \tangkapan{interpol2}
    \end{minipage}\hfill
    \begin{minipage}{0.48\textwidth}
        \centering
        \tangkapan{interpol2_eval1}
    \end{minipage}\hfill
    \begin{minipage}{0.48\textwidth}
        \centering
        \tangkapan{interpol2_eval2}
    \end{minipage}
\end{center}
\end{kasus}

\textbf{Analisis.} Berdasarkan pengujian, program berhasil memproses enam titik data awal yang menghasilkan persamaan polinomial $P(x) = 25 + 2.893x + 0.07x^2 + 0.02x^3 - 0.002x^4$. Persamaan  tersebut digunakan  untuk memprediksi peningkatan suhu secara maju menghasilkan taksiran berurutan sebesar $27.981^\circ\text{C}$ pada menit ke-1, $36.617^\circ\text{C}$ pada menit ke-3.5, $52.571^\circ\text{C}$ pada menit ke-7, dan $65.425^\circ\text{C}$ pada menit ke-9.2. Sementara itu, untuk memecahkan persoalan waktu kerusakan komponen pada suhu spesifik $65^\circ\text{C}$  program  menukar variabel suhu menjadi sumbu $x$ dan waktu menjadi sumbu $y$. Melalui pendekatan tersebut, terbentuk persamaan balikan baru $P(x) = -12.719 + 0.702x - 0.009x^2$ yang saat dievaluasi pada target $P(65)$ menunjukkan bahwa komponen akan rusak pada menit ke-9.14.

\subsection{Studi Kasus Interpolasi \textit{Natural Cubic Spline}}

\begin{kasus}
\ubag{Input}
\begin{center}
\begin{tabular}{c|ccccc}
$x_i$ & 1 & 2 & 3 & 4 & 5 \\ \hline
$y_i$ & 2 & 3 & 4 & 5 & 6
\end{tabular}
\end{center}
\ubag{Output}
Turunan kedua: $M_0 = M_1 = M_2 = M_3 = M_4 = 0$.
\begin{align*}
S_0(x) &= 2 + (x - 1), & 1 \le x \le 2, \\
S_1(x) &= 3 + (x - 2), & 2 \le x \le 3, \\
S_2(x) &= 4 + (x - 3), & 3 \le x \le 4, \\
S_3(x) &= 5 + (x - 4), & 4 \le x \le 5 .
\end{align*}
Contoh evaluasi: $S(2.5) = 3.5$.

\tangkapan{ncs_1}
\tangkapan{ncs_2}
\tangkapan{ncs_3}
\tangkapan{ncs_4}
\tangkapan{ncs_5}
\tangkapan{ncs_6}
\tangkapan{ncs_7}
\tangkapan{ncs_8}
\tangkapan{ncs_9}
\tangkapan{ncs_10}
\tangkapan{ncs_11}
\tangkapan{ncs_12}
\tangkapan{ncs_13}
\tangkapan{ncs_14}
\tangkapan{ncs_15}
\tangkapan{ncs_16}
\tangkapan{ncs_17}
\end{kasus}

\textbf{Analisis.} Data titik-titik pada kasus ini membentuk garis lurus $y = x + 1$. 
Hasil interpolasi menunjukkan bahwa seluruh turunan kedua bernilai nol, 
sehingga setiap segmen spline juga berbentuk garis lurus. 
Dengan demikian, keempat segmen menghasilkan fungsi yang sama, yaitu 
$y = x + 1$. Contohnya, pada $x = 2.5$ diperoleh $S(2.5) = 3.5$.

\subsection{Studi Kasus Regresi Spline}

\begin{kasus}
\ubag{Input}
\begin{center}
\begin{tabular}{c|ccccccc}
$i$   & 0 & 1 & 2 & 3 & 4 & 5 & 6 \\ \hline
$x_i$ & $-3$ & $-2$ & $-1$ & 0 & 1 & 2 & 3 \\
$y_i$ & 19.9 & 4.4 & $-0.5$ & 2 & 6.5 & 15.6 & 38.1
\end{tabular}
\end{center}
Spline kubik ($d = 3$) dengan satu knot internal $\xi_1 = 0$, sehingga jumlah basisnya $M = 3 + 1 + 1 = 5$ untuk $N = 7$ titik data. Basis yang dipakai adalah $1, x, x^2, x^3, (x - 0)_+^3$.
\ubag{Output}
\[
\hat y(x) = 2 + 3x - x^3 + 2(x)_+^3
\quad\Longleftrightarrow\quad
\hat y(x) =
\begin{cases}
2 + 3x - x^3, & x \le 0,\\
2 + 3x + x^3, & x > 0.
\end{cases}
\]
\begin{center}
\begin{tabular}{c|ccccccc}
$x_i$ & $-3$ & $-2$ & $-1$ & 0 & 1 & 2 & 3 \\ \hline
$y_i$ & 19.9 & 4.4 & $-0.5$ & 2 & 6.5 & 15.6 & 38.1 \\
$\hat y(x_i)$ & 20 & 4 & 0 & 2 & 6 & 16 & 38 \\
$e_i = y_i - \hat y(x_i)$ & $-0.1$ & $0.4$ & $-0.5$ & $0$ & $0.5$ & $-0.4$ & $0.1$
\end{tabular}
\end{center}
Jumlah kuadrat galat $\sum e_i^2 = 0.84$. Contoh evaluasi: $\hat y(-2.5) = 10.125$, $\hat y(0.5) = 3.625$, $\hat y(1.5) = 9.875$.
\tangkapan{regresi_1}
\tangkapan{regresi_2}
\tangkapan{regresi_3}
\tangkapan{regresi_4}
\tangkapan{regresi_5}
\tangkapan{regresi_6}
\tangkapan{regresi_7}
\tangkapan{regresi_8}
\end{kasus}

\textbf{Analisis.} Program membentuk matriks desain $X$ berukuran $7 \times 5$, lalu menyelesaikan persamaan normal dengan eliminasi Gauss. Sebelum itu, $x$ diubah menjadi $u = x/3$ karena data berada di $[-3, 3]$ (titik tengahnya 0 dan setengah lebarnya 3). Koefisien yang diperoleh dalam $u$ adalah $b = (2,\ 9,\ 0,\ -27,\ 54)$. Setelah dikembalikan ke $x$, yaitu dengan membagi $b_j$ dengan $3^j$ untuk suku polinom dan membagi $b_4$ dengan $3^3$ untuk suku \textit{knot}, didapat $\beta = (2,\ 3,\ 0,\ -1,\ 2)$. Untuk memastikannya, kami juga menyelesaikan persamaan normal dalam $x$ secara eksak dengan pecahan:
\[
X^TX =
\begin{bmatrix}
7 & 0 & 28 & 0 & 36\\
0 & 28 & 0 & 196 & 98\\
28 & 0 & 196 & 0 & 276\\
0 & 196 & 0 & 1588 & 794\\
36 & 98 & 276 & 794 & 794
\end{bmatrix},
\qquad
X^Ty =
\begin{bmatrix} 86\\ 84\\ 608\\ 588\\ 1160 \end{bmatrix},
\]
dan hasilnya sama persis.

Koefisien $\beta_4 = 2$ mengubah koefisien $x^3$ dari $-1$ di kiri \textit{knot} menjadi $-1 + 2 = 1$ di kanan \textit{knot}. Kedua potongan kurva tetap tersambung mulus di $x = 0$ karena nilainya sama-sama 2, turunan pertamanya ($3 - 3x^2$ dan $3 + 3x^2$) sama-sama 3, dan turunan keduanya ($-6x$ dan $6x$) sama-sama 0. Perubahan koefisien kubik inilah yang membuat bentuk kurva di kiri dan kanan \textit{knot} berbeda. Di sisi kiri, $\hat y$ turun dari sekitar 20 ke 0, lalu naik lagi, sedangkan di sisi kanan, $\hat y$ naik makin curam sampai 38.

Karena $N = 7$ lebih besar daripada $M = 5$, kurva tidak harus melewati semua titik. Galat terbesarnya hanya $0.5$, yaitu di $x = -1$ dan $x = 1$. Polanya juga simetris dengan tanda berlawanan, $e(-x) = -e(x)$, sehingga model tidak cenderung terlalu tinggi di satu sisi saja. Jumlah kuadrat galatnya $0.84$, sangat kecil dibandingkan dengan total variasi data yang sekitar $1100$yang dimana artinya, model menjelaskan lebih dari $99.9\%$ variasi $y$. Kami menduga data ini dibuat dari kurva $2 + 3x \mp x^3$ yang diberi sedikit gangguan, dan regresi berhasil mendapatkan kembali koefisien kurva tersebut.
